\documentclass[10pt,a4paper]{amsart}
\usepackage[margin=19mm]{geometry}
\usepackage{amsmath,amssymb,mathtools}
\usepackage[scr=rsfs,cal=euler]{mathalfa}
\usepackage{xfrac}
\usepackage[unicode,hidelinks,bookmarks=false]{hyperref}
\newcommand{\ie}{i.e.\ }
\newcommand{\cf}{cf.\ }
\newcommand{\resp}{respectively\ }
\newcommand{\Subsection}{\subsection}
\providecommand{\nobreakdash}{\nobreak\hbox{-}}
\setlength{\emergencystretch}{2em}
\multlinegap0pt
\usepackage{bm}
\usepackage{bbm}
\usepackage{dsfont}
\usepackage{xspace}
\usepackage{cancel}
\usepackage{mathtools}
\usepackage{cleveref}
\usepackage{amsthm}
\makeatletter
\@ifundefined{theorem}{\newtheorem{theorem}{Theorem}}{}
\@ifundefined{proposition}{\newtheorem{proposition}[theorem]{Proposition}}{}
\makeatother
\def \Mvec{\mathbf{M}}
\def\Hvec{\mathbf{H}}
\def\Qvec{\mathbf{Q}}
\title[A variational approach to ferronematics with a dimension red]{A variational approach to ferronematics with a dimension reduction}
\author{Shilpa Dutta}
\date{Source: arXiv:2606.01430v1 (CC BY 4.0)}
\begin{document}
\maketitle

\noindent\emph{Source and licence.} A variational approach to ferronematics with a dimension reduction. Version arXiv:2606.01430v1 (CC BY 4.0), distributed under the Creative Commons Attribution 4.0 International licence, as recorded in the arXiv licence metadata for this exact version and shown on the abstract page https://arxiv.org/abs/2606.01430v1. Full rights notice: licenses/arxiv-2606.01430-CC-BY-4.0.txt.\par\smallskip

\noindent\textit{Editorial scope.} Counted unit: the Proposition whose source label is \texttt{reduced-QM} in the article (source0.tex lines 714--722, source label \texttt{reduced-QM}), delivered together with the article's own complete original proof of it (source0.tex lines 723--763, one \texttt{proof} environment of 41 lines). No mathematics is written, repaired or re-derived: statement and proof are the article's own text line for line. Required context retained uncounted: bulk\_energy (lines 126--129); energy\_bound (lines 370--377); th:minimizer (lines 541--543). Every definition, display and label of those lines is retained unchanged. Omitted from this package and not redistributed: 1--125, 130--369, 378--540, 544--713, 764--958. Nothing inside a retained excerpt is omitted or altered: every retained body is delivered exactly as the article's lines are written, so no omission receipt of any kind accompanies this package. Citations: the retained excerpts carry no citation command at all, so no bibliography entry of the article is transcribed and no reference is redistributed. Reference adaptation: none was needed and none was made; every reference inside the delivered unit resolves inside it, so no unresolved reference or citation is produced. Printed slips kept literal: no printed word, symbol, hypothesis or number of the counted statement or of its proof was altered. Layout and class: the article's own class is \texttt{article} with packages including geometry, amsthm, titling, mathtools, xspace, mathrsfs, esint, algorithm2e, algorithmic, graphicx, tikz, xcolor, caption, subcaption; the wrapper is the lane's pinned \texttt{amsart} editorial wrapper at 10pt on A4, which restates the theorem-like environments the retained text needs and adds the article's own macro definitions that the retained text needs (\texttt{\textbackslash Hvec}, \texttt{\textbackslash Mvec}, \texttt{\textbackslash Qvec}), with extra wrapper packages (bm, bbm, dsfont, xspace, cancel, mathtools, cleveref). The delivered numbering is the unit's own. Assets: no retained span contains a figure environment, a graphics-inclusion command, a PDF-page inclusion, a table or a TikZ picture; no image file is copied into this package, the package declares no asset dependency and no image file is redistributed with it. Licence: version arXiv:2606.01430v1 of this article is distributed under the Creative Commons Attribution 4.0 International licence, as recorded in the arXiv licence metadata for this exact version and shown on the abstract page https://arxiv.org/abs/2606.01430v1. A full rights notice is delivered at licenses/arxiv-2606.01430-CC-BY-4.0.txt. Characters: every retained excerpt is pure ASCII, so the delivered engine, XeLaTeX with its default Latin Modern text font, reports no missing character.
\begin{eqnarray}\label{bulk_energy}
\mathcal{E}_B\left(\Qvec, \Mvec\right) = \int_{\Omega}\biggl[ \frac{K_1}{2}|\nabla \Qvec|^2 + \frac{K_2}{2}|\nabla \Mvec|^2 + f_B\left(\Qvec, \Mvec\right) + \frac{\mu_0}{2}|\Hvec_{\Mvec}|^2 - \mu_0\Mvec\cdot \Hvec_{ext}\nonumber\\
- \frac{\mu_0\gamma}{2}\Qvec\Mvec\cdot\Mvec - \frac{\chi_1\mu_0}{2}\Qvec\Hvec_{\Mvec}\cdot\Hvec_{\Mvec} - \frac{\chi_1\mu_0}{2}\Qvec\Hvec_{ext}\cdot\Hvec_{ext}\biggr]\;.
\end{eqnarray}

\begin{proposition}\label{energy_bound}
Let $\Omega\subset\mathbb R^3$ be a bounded domain such that $\partial\Omega$ is Lipschitz. Assume that $\left(\Qvec^{\varepsilon}, \Mvec^{\varepsilon}\right)$ be a sequence of minimizer of the functional $\mathcal{E}^{\varepsilon}_B$. Then, up to subsequences 
\begin{eqnarray}
\Qvec^{\varepsilon}\rightharpoonup\Tilde{\Qvec} &\mbox{ weakly in }& W^{1,2}\left(\Omega; \mathbb R^{3\times 3}\right),\\
\Mvec^{\varepsilon}\rightharpoonup\Tilde{\Mvec} &\mbox{ weakly in }& W^{1,2}\left(\Omega; \mathbb R^3\right).
\end{eqnarray}
Moreover, it can be shown that $\left(\Tilde{\Qvec}, \Tilde{\Mvec}\right)$ is a minimizer of the energy functional $\mathcal{E}_B\left(\Qvec, \Mvec\right)$.
\end{proposition}

\begin{theorem}\label{th:minimizer}
Let $\Omega_{\eta}\subset\mathbb R^3$ be a bounded domain such that $\mathcal{Q}_{\eta}\times\mathcal{M}_{\eta}\ne\emptyset$. Moreover, we assume that $\Hvec_{ext}\in C\left(\overline{\Omega}_{\eta}; \mathbb R^3\right)$. Then the functional $\mathcal{E}_{\eta}\left(\Qvec, \Mvec\right)$ (cf. \eqref{bulk_energy}) has a minimum on the space $\mathcal{Q}_{\eta}\times\mathcal{M}_{\eta}$.
\end{theorem}

\begin{proposition}\label{reduced-QM}
Let $\left(\Qvec_{\eta}, \Mvec_{\eta}\right)$ be a sequence of minimizers of the functional $\mathcal{E}_{\eta}\left(\Qvec, \Mvec\right)$. Then the following convergences holds, up to a subsequence
\begin{eqnarray}
\nabla_p\Qvec_{\eta}&\rightarrow&\nabla_p\Qvec \mbox{ strongly in } L^2\left(\Omega; \mathbb R^{3\times 3}\right),\label{eq:strong_gQ}\\
\frac{1}{\eta}\Qvec_{\eta,3}&\rightarrow& 0 \mbox{ strongly in } L^2\left(\Omega; \mathbb R^{3\times 3}\right),\label{eq:strong1_gQ}\\
\nabla_p\Mvec_{\eta}&\rightarrow&\nabla_p\Mvec \mbox{ strongly in } L^2\left(\Omega; \mathbb R^{3}\right),\label{eq:strong_gM}\\
\frac{1}{\eta}\Mvec_{\eta,3}&\rightarrow& 0 \mbox{ strongly in } L^2\left(\Omega; \mathbb R^{3}\right)\label{eq:strong1_gM}.
\end{eqnarray}
\end{proposition}

\begin{proof}
Thanks to \Cref{th:minimizer}, we assume that $\left(\Qvec_{\eta}, \Mvec_{\eta}\right)$ be a sequence of minimizers of the energy functional $\mathcal{E}_{\eta}\left(\Qvec, \Mvec\right)$, yields
\begin{eqnarray}
\mathcal{E}_{\eta}\left(\Qvec_{\eta}, \Mvec_{\eta}\right)\le\mathcal{E}_{\eta}\left(\Qvec, \Mvec\right)<\infty\label{eq:bound}
\end{eqnarray}
for some fixed $\left(\Qvec, \Mvec\right)\in \mathcal{A}_{\eta}\times\mathcal{S}_{\eta}$. Therefore, we have
\begin{eqnarray}
||\nabla_p\Mvec_{\eta}||_{L^2}\le D_1, \bigg|\bigg|\frac{1}{\eta}\Mvec_{\eta,3}\bigg|\bigg|_{L^2}\le D_2.\label{eq:M_conv1}
%\label{eq:M_conv2}
\end{eqnarray}
Thus, up to a subsequence,
\begin{eqnarray}
\nabla_p\Mvec_{\eta}&\rightharpoonup&\nabla_p\Mvec \mbox{ weakly in } L^2,\label{eq:M_conv3}\\
\frac{1}{\eta}\Mvec_{\eta,3}&\rightharpoonup& L_1  \mbox{ weakly in } L^2.\label{eq:M_conv4}
\end{eqnarray}
Similarly,
\begin{eqnarray}
||\nabla_p\Qvec_{\eta}||_{L^2}\le D'_1, \bigg|\bigg|\frac{1}{\eta}\Qvec_{\eta,3}\bigg|\bigg|_{L^2}\le D'_2.\label{eq:Q_conv2}
\end{eqnarray}
Therefore, up to a subsequence, we have
\begin{eqnarray}
\nabla_p\Qvec_{\eta}&\rightharpoonup&\nabla_p\Qvec \mbox{ weakly in } L^2,\label{eq:Q_conv3}\\
\frac{1}{\eta}\Qvec_{\eta,3}&\rightharpoonup& L_2  \mbox{ weakly in } L^2.\label{eq:Q_conv4}
\end{eqnarray}
We select $\Qvec_{\eta}=\Qvec$ such that $f^0_{sur}\left(\Qvec, \hat{x}_3\right)=0$ or $\Qvec\hat{x}_3=\beta\hat{x}_3$, and $\Mvec_{\eta}=\Mvec$ such that $\Mvec_{,3}=0$ as a recovery sequence, and obtain
\begin{eqnarray}
\limsup_{\eta\rightarrow 0}\int_{\Omega}|\nabla_p\Qvec_{\eta}|^2+\frac{1}{\eta^2}\Qvec_{\eta,3}+|\nabla_p\Mvec_{\eta}|^2+\frac{1}{\eta^2}\Mvec_{\eta,3}\le\int_{\Omega}|\nabla_p\Qvec|^2+|\nabla_p\Mvec|^2.\label{lim-sup-eq}
\end{eqnarray}
Applying the lower semicontinuity property of the $L^2$-norm to convergences \eqref{eq:M_conv3}, \eqref{eq:M_conv4}, \eqref{eq:Q_conv3} and \eqref{eq:Q_conv4}, and using the limsup inequality \eqref{lim-sup-eq}, we have
\begin{eqnarray}
& &\int_{\Omega}|\nabla_p\Qvec|^2+|\nabla_p\Mvec|^2+L^2_1+L_2^2\nonumber\\
&\le&\liminf_{\eta\rightarrow 0}\int_{\Omega}|\nabla_p\Qvec_{\eta}|^2+\frac{1}{\eta^2}\Qvec_{\eta,3}+|\nabla_p\Mvec_{\eta}|^2+\frac{1}{\eta^2}\Mvec_{\eta,3}\nonumber\\
&\le&\limsup_{\eta\rightarrow 0}\int_{\Omega}|\nabla_p\Qvec_{\eta}|^2+\frac{1}{\eta^2}\Qvec_{\eta,3}+|\nabla_p\Mvec_{\eta}|^2+\frac{1}{\eta^2}\Mvec_{\eta,3}\nonumber\\
&\underset{\eqref{lim-sup-eq}}{\le}&\int_{\Omega}|\nabla_p\Qvec|^2+|\nabla_p\Mvec|^2.\label{eq:limit_ineq}
\end{eqnarray}
By \eqref{eq:limit_ineq}, $L_1=L_2=0$ and
\begin{eqnarray}
\lim_{\eta\rightarrow 0}\int_{\Omega}|\nabla_p\Qvec_{\eta}|^2+\frac{1}{\eta^2}\Qvec_{\eta,3}+|\nabla_p\Mvec_{\eta}|^2+\frac{1}{\eta^2}\Mvec_{\eta,3}=\int_{\Omega}|\nabla_p\Qvec|^2+|\nabla_p\Mvec|^2.
\end{eqnarray} 
The strong convergences \eqref{eq:strong_gQ}, \eqref{eq:strong1_gQ}, \eqref{eq:strong_gM} and \eqref{eq:strong1_gM} follow by repeating the same approach as in Proposition \ref{energy_bound}.
\end{proof}

\end{document}
